\documentclass[UTF8,a4paper,12pt]{ctexart}

\usepackage{amsmath,amssymb}
\usepackage{booktabs}
\usepackage{caption}
\usepackage{float}
\usepackage{geometry}
\usepackage{graphicx}
\usepackage{hyperref}
\usepackage{longtable}
\usepackage{xcolor}
\usepackage{enumitem}
\usepackage{fvextra}

\geometry{left=2.6cm,right=2.6cm,top=2.6cm,bottom=2.6cm}
\hypersetup{colorlinks=true,linkcolor=black,urlcolor=blue}
\setlength{\parindent}{2em}
\setlength{\parskip}{0.25em}
\setlength{\emergencystretch}{3em}
\linespread{1.25}
\setlist[itemize]{leftmargin=2.5em,itemsep=0.15em}
\DefineVerbatimEnvironment{CodeBlock}{Verbatim}{breaklines,breakanywhere,fontsize=\small,frame=single,rulecolor=\color{gray!45}}

\title{\heiti 机器学习与Python实践\\[0.8em]\Large 第1次作业\\[1.6em]\Huge 使用线性回归构建房价预测模型}
\author{}
\date{}

\begin{document}

\begin{titlepage}
    \centering
    \vspace*{1.2cm}
    
\includegraphics[width=0.78\linewidth]{figures/cover_logo.png}\par
    \vspace{1.0cm}
    {\zihao{2}\heiti 机器学习与Python实践\par}
    \vspace{0.8cm}
    {\zihao{3}\heiti 第1次作业\par}
    \vspace{1.2cm}
    {\zihao{1}\heiti 使用线性回归构建房价预测模型\par}
    \vfill
    \begin{tabular}{rl}
        姓名： & 匿名\\[0.45em]
        专业： & 强基班\\[0.45em]
        学号： & \\[0.45em]
        教师： & 匿名\\[0.45em]
        学院： & 信息科学与工程学院\\[0.45em]
    \end{tabular}
    \vfill
    {\zihao{4}2026年05月08日\par}
\end{titlepage}

\tableofcontents
\newpage

\section{前言及问题概述}

线性回归是最简单、最经典的回归模型之一。假设输入样本 $x$ 为 $d$ 维特征向量，预测目标 $y$ 为连续型取值，线性回归模型可写为
\begin{equation}
    \hat{y}=w_0+w_1x_1+\cdots+w_dx_d .
\end{equation}
其中 $w_0$ 为截距项，$w_1,\ldots,w_d$ 为各特征对应的回归系数。模型训练的核心任务是根据训练数据估计参数，使预测值与真实值之间的误差尽可能小。

本次实验围绕波士顿房价数据集构建房价预测模型。数据集中每个样本对应一个地区的住房及环境指标，目标变量为房屋价格中位数。实验使用 Python 及 \texttt{sklearn.linear\_model} 模块中的 \texttt{LinearRegression} 函数建立最小二乘线性回归模型，同时结合自定义批量梯度下降模型、随机梯度下降模型、岭回归、Lasso 回归和二次多项式回归进行对比。通过均方误差（MSE）和决定系数 $R^2$ 评价模型效果，并分析归一化、测试集比例和学习率等因素对模型性能的影响。

\section{相关基础知识及方法}

本章从线性回归原理、数据集字段含义和算法实现流程三个方面说明实验方法。报告正文遵循模板文件中的章节目录，具体实验内容依据原“报告1”中的代码、运行结果和分析文字整理。

\subsection{相关原理}

线性回归假设目标变量与输入特征之间存在近似线性关系。设训练集为 $\{(x^{(i)},y^{(i)})\}_{i=1}^{n}$，其中 $x^{(i)}\in \mathbb{R}^d$，加入截距项后可记为矩阵形式
\begin{equation}
    \hat{\mathbf{y}} = \mathbf{X}\mathbf{w}.
\end{equation}
最小二乘法通过最小化残差平方和估计模型参数：
\begin{equation}
    J(\mathbf{w})=\frac{1}{n}\sum_{i=1}^{n}(\hat{y}^{(i)}-y^{(i)})^2 .
\end{equation}
当 $\mathbf{X}^{T}\mathbf{X}$ 可逆时，正规方程可直接求得闭式解：
\begin{equation}
    \mathbf{w}=(\mathbf{X}^{T}\mathbf{X})^{-1}\mathbf{X}^{T}\mathbf{y}.
\end{equation}
在 \texttt{sklearn.linear\_model.LinearRegression} 中，回归系数即由最小二乘思想进行估计。与之相比，梯度下降方法不直接求闭式解，而是沿损失函数负梯度方向迭代更新参数：
\begin{equation}
    \mathbf{w}\leftarrow \mathbf{w}-\eta \nabla J(\mathbf{w}),
\end{equation}
其中 $\eta$ 为学习率。批量梯度下降每轮使用全部训练样本计算梯度，随机梯度下降每次或每小批次使用部分样本更新参数，适合样本量较大或在线学习场景。

本实验采用两个常用评价指标。均方误差 MSE 越小，表示预测值与真实值越接近；决定系数 $R^2$ 越接近 1，表示模型对目标变量波动的解释能力越强。

\subsection{数据描述}

波士顿房价数据集共包含 506 条样本，特征数为 13，目标变量为 \texttt{MEDV}，表示自住房屋价格中位数。实验代码从本地 \texttt{boston\_house\_prices.csv} 文件读取数据，并要求目标列名为 \texttt{MEDV}。主要特征含义见表 \ref{tab:features}。

\begin{longtable}{lll}
\caption{波士顿房价数据集主要字段说明}\label{tab:features}\\
\toprule
字段 & 类型 & 含义\\
\midrule
\endfirsthead
\toprule
字段 & 类型 & 含义\\
\midrule
\endhead
\texttt{CRIM} & 连续特征 & 城镇人均犯罪率\\
\texttt{ZN} & 连续特征 & 大面积住宅用地比例\\
\texttt{INDUS} & 连续特征 & 非零售商业用地比例\\
\texttt{CHAS} & 二值特征 & 是否邻近查尔斯河\\
\texttt{NOX} & 连续特征 & 一氧化氮浓度\\
\texttt{RM} & 连续特征 & 每户平均房间数\\
\texttt{AGE} & 连续特征 & 1940 年以前建成住房比例\\
\texttt{DIS} & 连续特征 & 到就业中心的加权距离\\
\texttt{RAD} & 连续特征 & 到放射状公路的可达性指数\\
\texttt{TAX} & 连续特征 & 房产税率\\
\texttt{PTRATIO} & 连续特征 & 城镇师生比例\\
\texttt{B} & 连续特征 & 人口结构相关指标\\
\texttt{LSTAT} & 连续特征 & 低收入人口比例\\
\texttt{MEDV} & 目标变量 & 自住房屋价格中位数\\
\bottomrule
\end{longtable}

由于各特征量纲差异较大，例如犯罪率、税率、房间数和距离的数值范围并不一致，因此实验在训练集上计算均值和标准差，对训练集及测试集进行 Z-score 标准化。该处理对正规方程求解影响较小，但对梯度下降类方法的收敛速度和稳定性影响明显。

\subsection{算法实现}

实验流程如下：
\begin{enumerate}
    \item 使用 \texttt{pandas} 读取 \texttt{boston\_house\_prices.csv}，将 \texttt{MEDV} 作为预测目标，其余列作为输入特征。
    \item 使用 \texttt{train\_test\_split} 按不同测试集比例划分训练集和测试集，主要比较 \texttt{test\_rate=0.45} 与 \texttt{test\_rate=0.2}。
    \item 对训练集计算均值和标准差，并用同一组统计量标准化训练集和测试集。
    \item 分别训练自定义批量梯度下降模型、\texttt{LinearRegression} 正规方程模型和 \texttt{SGDRegressor} 随机梯度下降模型。
    \item 输出训练时长、训练集 MSE、测试集 MSE、训练集 $R^2$ 与测试集 $R^2$。
    \item 绘制损失曲线、真实值与预测值散点图，并进一步测试学习率、正则化模型和多项式特征对预测效果的影响。
\end{enumerate}

核心实现中，自定义模型的前向传播为
\begin{equation}
    \hat{y} = Xw+b,
\end{equation}
梯度为
\begin{equation}
    \frac{\partial J}{\partial w}=\frac{1}{n}X^T(\hat{y}-y),\qquad
    \frac{\partial J}{\partial b}=\frac{1}{n}\sum_{i=1}^{n}(\hat{y}^{(i)}-y^{(i)}).
\end{equation}
随后按学习率进行参数更新。完整 Python 程序见附录。

\section{实验结果与可扩展内容}

\subsection{实验结果}

在 \texttt{test\_rate=0.45} 且进行标准化的情况下，三种线性模型均能较好拟合房价数据，测试集 $R^2$ 约为 0.73--0.74。具体结果见表 \ref{tab:main-result}。

\begin{table}[H]
\centering
\caption{标准化后 \texttt{test\_rate=0.45} 的模型结果}
\label{tab:main-result}
\begin{tabular}{lccccc}
\toprule
模型 & 训练时长/s & 训练 MSE & 训练 $R^2$ & 测试 MSE & 测试 $R^2$\\
\midrule
自定义批量梯度下降 & 0.0128 & 22.93 & 0.73 & 22.38 & 0.73\\
\texttt{LinearRegression} & 0.0101 & 22.60 & 0.74 & 21.87 & 0.74\\
\texttt{SGDRegressor} & 0.0011 & 22.67 & 0.73 & 22.13 & 0.73\\
\bottomrule
\end{tabular}
\end{table}

\begin{figure}[H]
    \centering
    \includegraphics[width=0.92\linewidth]{figures/original_02.png}
    \caption{标准化后损失曲线及三种模型预测效果}
    \label{fig:normalized}
\end{figure}

从图 \ref{fig:normalized} 可以看到，自定义批量梯度下降模型的损失值在训练前期快速下降，随后逐渐趋于稳定；三个模型的预测散点大体围绕理想预测线分布，说明线性模型能够捕捉房价与主要特征之间的基本关系。其中 \texttt{LinearRegression} 的测试集 MSE 最低，$R^2$ 最高，作为最小二乘闭式求解模型具有稳定优势。

将标准化处理去除后，正规方程模型的结果几乎不变，而随机梯度下降模型训练时间明显增加，测试效果略有下降，见表 \ref{tab:unnormalized}。

\begin{table}[H]
\centering
\caption{未标准化时的模型结果}
\label{tab:unnormalized}
\begin{tabular}{lccc}
\toprule
模型 & 训练时长/s & 测试 MSE & 测试 $R^2$\\
\midrule
\texttt{LinearRegression} & 0.0008 & 21.87 & 0.74\\
\texttt{SGDRegressor} & 1.7360 & 23.24 & 0.72\\
\bottomrule
\end{tabular}
\end{table}

\begin{figure}[H]
    \centering
    \includegraphics[width=0.82\linewidth]{figures/original_03.png}
    \caption{未标准化时正规方程与随机梯度下降预测效果}
    \label{fig:unnormalized}
\end{figure}

进一步改变学习率为 0.1 时，随机梯度下降出现明显不收敛现象，测试集 MSE 变得极大，$R^2$ 为大幅负值。这说明梯度下降对学习率敏感，过大的学习率会导致参数更新越过最优区域，造成震荡甚至发散。

\begin{figure}[H]
    \centering
    \includegraphics[width=0.58\linewidth]{figures/original_04.png}
    \caption{学习率过大时随机梯度下降预测发散}
    \label{fig:eta}
\end{figure}

当测试集比例设为 0.2 时，训练集 MSE 略低，但测试集 MSE 增大，测试集 $R^2$ 降至约 0.65--0.67，见表 \ref{tab:test02}。本次划分下，\texttt{test\_rate=0.45} 的测试表现更好。

\begin{table}[H]
\centering
\caption{标准化后 \texttt{test\_rate=0.2} 的模型结果}
\label{tab:test02}
\begin{tabular}{lccccc}
\toprule
模型 & 训练时长/s & 训练 MSE & 训练 $R^2$ & 测试 MSE & 测试 $R^2$\\
\midrule
自定义批量梯度下降 & 0.0138 & 21.95 & 0.75 & 25.64 & 0.65\\
\texttt{LinearRegression} & 0.0005 & 21.64 & 0.75 & 24.29 & 0.67\\
\texttt{SGDRegressor} & 0.0010 & 21.72 & 0.75 & 24.88 & 0.66\\
\bottomrule
\end{tabular}
\end{table}

\begin{figure}[H]
    \centering
    \includegraphics[width=0.92\linewidth]{figures/original_05.png}
    \caption{\texttt{test\_rate=0.2} 时损失曲线及预测效果}
    \label{fig:test02}
\end{figure}

\subsection{参数调优及可扩展内容}

在基础线性回归之外，报告 1 中还加入了岭回归、Lasso 回归和二次多项式回归进行扩展实验，结果见表 \ref{tab:extension}。

\begin{table}[H]
\centering
\caption{扩展模型测试结果}
\label{tab:extension}
\begin{tabular}{lccc}
\toprule
模型 & 训练时长/s & 测试 MSE & 测试 $R^2$\\
\midrule
岭回归 & 0.0014 & 24.31 & 0.67\\
Lasso 回归 & 0.0006 & 25.66 & 0.65\\
二次多项式回归 & 0.0030 & 14.29 & 0.81\\
\bottomrule
\end{tabular}
\end{table}

\begin{figure}[H]
    \centering
    \includegraphics[width=0.9\linewidth]{figures/original_06.png}
    \caption{岭回归、Lasso 回归与二次多项式回归预测效果}
    \label{fig:extension}
\end{figure}

结果显示，在本次实验设定下，二次多项式回归通过引入特征之间的非线性组合，测试集 MSE 降至 14.29，$R^2$ 提升到 0.81，优于基础线性模型。但多项式特征会增加模型复杂度，若阶数继续升高，可能出现过拟合。因此模型扩展时应结合交叉验证、正则化和测试集表现综合判断。

此外，报告 1 对 BGD、SGD 和 MBGD 三种梯度下降策略进行了比较。测试集结果为：BGD 的 MSE 为 16.8523，$R^2$ 为 0.7687；SGD 的 MSE 为 16.0348，$R^2$ 为 0.7799；MBGD 的 MSE 为 16.7107，$R^2$ 为 0.7707。该结果表明，小批量或随机更新在合适学习率下也能取得较好效果。

\begin{figure}[H]
    \centering
    \includegraphics[width=0.9\linewidth]{figures/original_08.png}
    \caption{BGD、SGD 与 MBGD 的预测及指标对比}
    \label{fig:gdcompare}
\end{figure}

\section{讨论与结论}

\subsection{结果讨论}

通过实验可以得到以下结论。第一，最小二乘线性回归在波士顿房价预测任务中可以作为有效基线模型。在 \texttt{test\_rate=0.45} 且标准化的条件下，\texttt{LinearRegression} 测试集 MSE 为 21.87，测试集 $R^2$ 为 0.74，说明模型能解释约 74\% 的房价波动。

第二，正规方程和梯度下降在理论上都可以求解线性回归参数，但计算方式不同。正规方程直接求闭式解，适合特征维度不高、样本规模适中的问题；梯度下降通过迭代逼近最优解，更适合大规模数据或需要在线更新的任务。本实验数据规模较小，因此 \texttt{LinearRegression} 的速度和稳定性表现较好。

第三，标准化对梯度下降影响明显。未标准化时，不同特征尺度差异会使损失函数等高线变得狭长，导致梯度下降更新方向不稳定，必须使用更小学习率才能收敛；而标准化使各特征处于相近尺度，优化路径更加平滑。

第四，学习率是梯度下降中的关键超参数。学习率过小会导致收敛缓慢，学习率过大则会导致震荡或发散。本实验中将随机梯度下降学习率调到 0.1 后，模型预测结果明显异常，验证了学习率调节的重要性。

\subsection{结论总结}

本次实验完成了基于波士顿房价数据集的线性回归建模与预测。实验先使用最小二乘线性回归建立房价预测模型，再与自定义批量梯度下降、随机梯度下降及若干扩展模型进行对比。结果表明，基础线性回归具有实现简单、速度快、解释性强等优点，适合作为回归任务的入门模型和基线模型；在存在非线性关系时，多项式特征可以进一步提升预测效果，但也需要注意控制模型复杂度。

综合来看，线性回归虽然形式简单，但完整实验过程包含数据读取、训练测试划分、特征标准化、模型训练、指标评价和可视化分析等机器学习基本环节。通过该案例可以较系统地理解监督学习回归任务的基本流程。

\section{心得体会与思考}

\subsection{个人心得}

通过本次实验，我对线性回归模型的理解从公式推导进一步落实到代码实现。最小二乘法可以直接给出回归系数，而梯度下降则通过不断迭代更新参数逐步减小误差。两种方法得到的模型形式相同，但适用场景和对数据预处理的敏感程度不同。

在编写程序和整理结果时，我也意识到评价模型不能只看训练集表现。训练集误差低并不代表模型泛化能力强，测试集 MSE 和 $R^2$ 更能反映模型对未知样本的预测效果。可视化散点图也很有帮助，它能直观展示预测点是否围绕理想直线分布。

\subsection{启发与思考}

线性回归模型的优势在于结构清晰、解释性强，但它对特征与目标之间的关系形式有较强假设。当真实关系存在明显非线性时，仅靠线性模型可能难以充分拟合。后续可以从三个方向继续改进：一是进行特征工程，构造更有意义的组合特征；二是使用交叉验证选择更稳定的测试结果；三是尝试岭回归、Lasso、多项式回归或树模型等方法，在模型复杂度和泛化能力之间取得平衡。

\section{附录}

\subsection*{算法代码及其注释}
\addcontentsline{toc}{subsection}{算法代码及其注释}

\VerbatimInput[breaklines,breakanywhere,fontsize=\small,frame=single,rulecolor=\color{gray!45}]{report_code.py}

\subsection*{前三次参数更新过程}
\addcontentsline{toc}{subsection}{前三次参数更新过程}

报告 1 中记录的前三次批量梯度下降参数更新如下：

\begin{CodeBlock}
迭代 1:
梯度向量:
[  3.5548 -3.2920  4.4314 -1.6078  3.9072 -6.3888  3.4429 -2.2753  3.4938  4.2933  4.6557 -3.0425  6.7698 -22.5340]
权重变化向量:
[ -0.0355  0.0329 -0.0443  0.0161 -0.0391  0.0639 -0.0344  0.0228 -0.0349 -0.0429 -0.0466  0.0304 -0.0677  0.2253]
当前权重向量:
[ -0.0409  0.0416 -0.0346  0.0195 -0.0490  0.0587 -0.0474  0.0319 -0.0332 -0.0354 -0.0525  0.0489 -0.0495  0.2241]
当前损失值: 592.1959

迭代 2:
梯度向量:
[  3.3399 -3.0786  4.1422 -1.5844  3.6319 -6.1818  3.1906 -2.0254  3.2288  4.0095  4.4563 -2.8697  6.4819 -22.3087]
权重变化向量:
[ -0.0334  0.0308 -0.0414  0.0158 -0.0363  0.0618 -0.0319  0.0203 -0.0323 -0.0401 -0.0446  0.0287 -0.0648  0.2231]
当前权重向量:
[ -0.0743  0.0724 -0.0760  0.0353 -0.0854  0.1205 -0.0793  0.0522 -0.0655 -0.0755 -0.0970  0.0776 -0.1143  0.4472]
当前损失值: 577.6884

迭代 3:
梯度向量:
[  3.1382 -2.8785  3.8709 -1.5611  3.3742 -5.9852  2.9544 -1.7919  2.9804  3.7432  4.2678 -2.7075  6.2104 -22.0856]
权重变化向量:
[ -0.0314  0.0288 -0.0387  0.0156 -0.0337  0.0599 -0.0295  0.0179 -0.0298 -0.0374 -0.0427  0.0271 -0.0621  0.2209]
当前权重向量:
[ -0.1057  0.1012 -0.1147  0.0509 -0.1191  0.1803 -0.1089  0.0701 -0.0953 -0.1129 -0.1397  0.1046 -0.1765  0.6680]
当前损失值: 563.8403
\end{CodeBlock}

\end{document}
